% In this file you should put all LaTeX macros and settings to be used both by
% the pdf version and the web version.
% This should be most of your macros.

% The theorem-like environments defined below are those that appear by default
% in the dependency graph. See the README of leanblueprint if you need help to
% customize this. 
% The configuration below use the theorem counter for all those environments
% (this is what the [theorem] arguments mean) and never resets it.
% If you want for instance to number them within chapters then you can add
% [chapter] at the end of the next line.
% All node environments are upright (definition style): their bodies are schematic -- lists,
% run-in labels, badges -- which read badly in the italics of the plain style.
\theoremstyle{definition}
\newtheorem{theorem}{Theorem}
\newtheorem{proposition}[theorem]{Proposition}
\newtheorem{lemma}[theorem]{Lemma}
\newtheorem{corollary}[theorem]{Corollary}

\theoremstyle{definition}
\newtheorem{definition}[theorem]{Definition}

% Assumption: something the development assumes rather than proves
% (a hypothesis bundle field, a standalone hypothesis binder, or an
% inherited trust boundary). Upright-titled like a theorem, so restore
% \theoremstyle{plain} after the definition-style block above.
\theoremstyle{definition}
\newtheorem{assumption}[theorem]{Assumption}

% Remark: prose asides that must NOT be dependency-graph nodes. Kept
% out of the thms= list in web.tex, and starred so it is unnumbered;
% plasTeX's amsthm shim explicitly supports the starred \newtheorem.
\theoremstyle{remark}
\newtheorem*{remark}{Remark}

% Notation for the lambda-box target language and its semantics
% (conventions.md sec. 7 -- writers use only these plus standard LaTeX).
\newcommand{\lbox}{\ensuremath{\lambda_\square}}
\newcommand{\lboxT}{\ensuremath{\lambda_\square^{T}}}
\newcommand{\Erases}{\ensuremath{\mathsf{Erases}}}
\newcommand{\ErasesEnv}{\ensuremath{\mathsf{ErasesEnv}}}
\newcommand{\Lower}{\ensuremath{\mathsf{Lower}}}
\newcommand{\LowerEnv}{\ensuremath{\mathsf{LowerEnv}}}
\newcommand{\SEval}{\ensuremath{\mathsf{SEval}}}
\newcommand{\WcbvEval}{\ensuremath{\mathsf{WcbvEval}}}
\newcommand{\TrExprS}{\ensuremath{\mathsf{TrExprS}}}
\newcommand{\Erasable}{\ensuremath{\mathsf{Erasable}}}
\newcommand{\mkApps}{\ensuremath{\mathsf{mkApps}}}
\newcommand{\code}[1]{\texttt{#1}}
\newcommand{\ie}{i.e.,}
\newcommand{\eg}{e.g.,}

% --- schematic-style helpers -------------------------------------------------
% Run-in label opening a line of a node: \lead{In short}, \lead{Given}, \lead{Then},
% \lead{Lean}, \lead{Idea}, \lead{Steps}, \lead{Caveat}, \lead{Deviation}, \lead{Status}.
\newcommand{\lead}[1]{\textbf{#1.}}
% Highlight of the one phrase a skimming reader must not miss.
\definecolor{bphl}{HTML}{9A3B00}
\newcommand{\hl}[1]{\textcolor{bphl}{\textbf{#1}}}
% Status badges. Use exactly one where a status is asserted.
\definecolor{bpproved}{HTML}{1B7F3B}
\definecolor{bpchecked}{HTML}{0B6E99}
\definecolor{bpassumed}{HTML}{B35900}
\definecolor{bpinherited}{HTML}{6A3FA0}
\definecolor{bpopen}{HTML}{B00020}
\definecolor{bpoutside}{HTML}{555555}
\newcommand{\stProved}{\textcolor{bpproved}{\textbf{[PROVED]}}}           % proved in this repository
\newcommand{\stChecked}{\textcolor{bpchecked}{\textbf{[CHECKED]}}}        % discharged by a checked term / kernel computation
\newcommand{\stAssumed}{\textcolor{bpassumed}{\textbf{[ASSUMED]}}}        % a hypothesis binder or bundle field
\newcommand{\stInherited}{\textcolor{bpinherited}{\textbf{[INHERITED]}}}  % trust inherited from lean4lean / Lean axioms
\newcommand{\stOpen}{\textcolor{bpopen}{\textbf{[OPEN]}}}                 % open proof obligation
\newcommand{\stOutside}{\textcolor{bpoutside}{\textbf{[OUTSIDE]}}}        % outside the verified perimeter

% --- Chapter 13/1 helpers: source tags, divergence-kind badges, divergence ids ----
% One colour per cited source. \src{JACM} -> coloured bold [JACM]; \srcs{JACM}{Sec. 7.3}
% -> coloured bold [JACM, Sec. 7.3]. Used only in chapters 1 and 13.
\definecolor{srcJACM}{HTML}{0072B2}
\definecolor{srcCCC}{HTML}{56B4E9}
\definecolor{srcHAB}{HTML}{009E73}
\definecolor{srcMR}{HTML}{999999}
\definecolor{srcMRD}{HTML}{CC79A7}
\definecolor{srcL}{HTML}{D55E00}
\definecolor{srcR}{HTML}{E69F00}
\definecolor{srcCT}{HTML}{6A3FA0}
\newcommand{\src}[1]{\textcolor{src#1}{\textbf{[#1]}}}
\newcommand{\srcs}[2]{\textcolor{src#1}{\textbf{[#1, #2]}}}

% Divergence-kind badges (chapter 13's six kinds). Three FORCED-* share one hue
% (reds/oranges) shading from lightest (Lean) to darkest (lean4lean).
\definecolor{kindFL}{HTML}{CC3311}
\definecolor{kindFE}{HTML}{E67300}
\definecolor{kindFK}{HTML}{9E4300}
\definecolor{kindDS}{HTML}{0072B2}
\definecolor{kindOP}{HTML}{D6006D}
\definecolor{kindGP}{HTML}{7F7F7F}
\newcommand{\kFL}{\textcolor{kindFL}{\textbf{[F-Lean]}}}
\newcommand{\kFE}{\textcolor{kindFE}{\textbf{[F-Eraser]}}}
\newcommand{\kFK}{\textcolor{kindFK}{\textbf{[F-l4l]}}}
\newcommand{\kDS}{\textcolor{kindDS}{\textbf{[Design]}}}
\newcommand{\kOP}{\textcolor{kindOP}{\textbf{[Open]}}}
\newcommand{\kGP}{\textcolor{kindGP}{\textbf{[Gap]}}}

% Divergence ids (\did{T1}, \did{B2}, ...): bold, coloured by the area letter
% prefix (T target calculus, B erasability, R erasure relation, V environment,
% C correctness, P pipeline, M method, U trust, S scope).
\definecolor{didT}{HTML}{0072B2}
\definecolor{didB}{HTML}{009E73}
\definecolor{didR}{HTML}{CC3311}
\definecolor{didV}{HTML}{E67300}
\definecolor{didC}{HTML}{6A3FA0}
\definecolor{didP}{HTML}{CC79A7}
\definecolor{didM}{HTML}{D6006D}
\definecolor{didU}{HTML}{555555}
\definecolor{didS}{HTML}{228833}
\newcommand{\didSetColor}[1]{%
  \ifx#1T\def\didColorName{didT}\else
  \ifx#1B\def\didColorName{didB}\else
  \ifx#1R\def\didColorName{didR}\else
  \ifx#1V\def\didColorName{didV}\else
  \ifx#1C\def\didColorName{didC}\else
  \ifx#1P\def\didColorName{didP}\else
  \ifx#1M\def\didColorName{didM}\else
  \ifx#1U\def\didColorName{didU}\else
  \ifx#1S\def\didColorName{didS}\else
  \def\didColorName{didU}%
  \fi\fi\fi\fi\fi\fi\fi\fi\fi}
\newcommand{\did}[1]{\didDispatch#1\didStop{#1}}
\def\didDispatch#1#2\didStop#3{\didSetColor{#1}\textcolor{\didColorName}{\textbf{#3}}}
